\lecture{45}{Неоднородность клиентов}

\sourcevideo
    {https://www.youtube.com/playlist?list=PLOzRYVm0a65fhKDS8cYIC7YCuVGJ4FOJx}
    {Week 12: Lecture 45}

В федеративном обучении выбранные клиенты начинают локальную работу
из одной серверной модели, но могут выполнить различный объём
вычислений. Причинами служат неодинаковые размеры локальных выборок,
скорости устройств, временные бюджеты, шаги и оптимизаторы. В
результате сервер усредняет обновления, полученные по различным
локальным траекториям.

\section{Локальная работа в FedAvg}

Пусть на коммуникационном раунде \(t\) сервер хранит модель
\(x^t\in\mathbb R^d\) и выбирает множество клиентов \(S_t\). Клиент
\(i\in S_t\) инициализирует
\[
    x_i^{t,0}=x^t
\]
и выполняет \(\tau_i\) локальных стохастических шагов
\[
    x_i^{t,j+1}
    =
    x_i^{t,j}
    -
    \eta_i g_i
    \left(
        x_i^{t,j};\xi_i^{t,j}
    \right),
    \qquad
    j=0,\ldots,\tau_i-1.
\]
После этого сервер агрегирует полученные модели:
\[
    x^{t+1}
    =
    \sum_{i\in S_t}p_i^t x_i^{t,\tau_i},
    \qquad
    \sum_{i\in S_t}p_i^t=1.
\]
При взвешивании по числу объектов обычно берут
\[
    p_i^t
    =
    \frac{n_i}{\sum_{r\in S_t}n_r},
\]
где \(n_i\) --- размер локальной выборки.

Если клиент использует мини-пакет размера \(B_i\) и выполняет
\(E_i\) эпох, то число локальных обновлений равно
\[
    \tau_i
    =
    E_i
    \left\lceil\frac{n_i}{B_i}\right\rceil.
\]
При делимости \(n_i\) на \(B_i\) получаем
\(\tau_i=E_in_i/B_i\). Таким образом, одинаковое число эпох не
означает одинакового числа шагов.

\section{Виды неоднородности}

Статистическая неоднородность означает различие локальных
распределений и функций:
\[
    \mathcal D_i\ne\mathcal D_j,
    \qquad
    F_i\ne F_j.
\]
Вычислительная неоднородность относится к различиям в самом процессе
локального обучения: значениях \(\tau_i\), \(\eta_i\), времени одного
шага и типе оптимизатора. Эти эффекты независимы: одинаковые
устройства могут хранить выборки разного размера, а одинаковые по
распределению данные могут обрабатываться с разной скоростью.

При общих \(E\) и \(B\)
\[
    n_i>n_j
    \quad\Longrightarrow\quad
    \tau_i\ge\tau_j.
\]
Локальное изменение модели имеет вид
\[
    \Delta_i^t
    :=
    x_i^{t,\tau_i}-x^t
    =
    -\eta_i
    \sum_{j=0}^{\tau_i-1}
    g_i
    \left(
        x_i^{t,j};\xi_i^{t,j}
    \right).
\]
Поэтому размер выборки влияет на глобальное обновление дважды: через
вес \(p_i^t\) и через длину локальной траектории \(\tau_i\). Первый
эффект отражает статистический вес клиента, второй изменяет само
направление и масштаб передаваемого обновления.

\section{Скорость устройств и отстающие клиенты}

Пусть \(c_i>0\) --- время одного локального шага клиента \(i\). При
фиксированном числе шагов \(\tau\) его локальное время примерно равно
\[
    T_i^{\mathrm{loc}}=\tau c_i.
\]
В синхронной схеме сервер начинает следующий раунд только после
получения требуемых ответов, поэтому
\[
    T_t^{\mathrm{round}}
    \approx
    \max_{i\in S_t}T_i^{\mathrm{loc}}
    +
    T_t^{\mathrm{comm}}.
\]
Клиенты с большим \(c_i\) становятся отстающими и определяют
продолжительность раунда.

Фиксация общего числа шагов выравнивает алгоритмический объём
локальной работы, но не время завершения. Кроме того, равное
\(\tau\) не означает одинакового числа просмотренных эпох, если
\(n_i\) различаются.

\section{Фиксированный временной бюджет}

Другой вариант состоит в выделении каждому клиенту одинакового
времени \(H>0\). Тогда
\[
    \tau_i
    =
    \left\lfloor\frac{H}{c_i}\right\rfloor,
\]
и более быстрый клиент выполняет больше шагов. Продолжительность
локальной стадии становится контролируемой, но сервер получает
обновления, построенные по траекториям различной длины.

Тем самым при разных скоростях устройств нельзя одновременно
обеспечить одинаковое число локальных шагов и одинаковое время
завершения. Выбор между этими режимами определяется тем, важнее ли
синхронизировать объём работы или ограничить задержку раунда.

\section{Различия шагов и локальных оптимизаторов}

Даже при одинаковом \(\tau_i\) масштаб \(\Delta_i^t\) зависит от
шага \(\eta_i\). Слишком большой шаг усиливает локальный дрейф и
может нарушить устойчивость, а слишком малый делает вклад клиента
незначительным.

Клиенты могут также использовать различные оптимизаторы. Например,
локальный метод с импульсом имеет вид
\[
    v_i^{t,j+1}
    =
    \beta_i v_i^{t,j}
    +
    g_i
    \left(
        x_i^{t,j};\xi_i^{t,j}
    \right),
\]
\[
    x_i^{t,j+1}
    =
    x_i^{t,j}
    -
    \eta_i v_i^{t,j+1}.
\]
Различия в \(\eta_i\), \(\beta_i\), состоянии импульса и типе
оптимизатора означают, что одно и то же число шагов не задаёт один и
тот же эффективный объём локального изменения. Сервер, получающий
только конечную модель, не всегда может восстановить эту информацию.

\section{Нормализация обновлений}

Обычная агрегация через локальные приращения записывается как
\[
    x^{t+1}
    =
    x^t
    +
    \sum_{i\in S_t}p_i^t\Delta_i^t.
\]
При различных \(\tau_i\) она смешивает обновления разного масштаба.
Естественная коррекция состоит в нормализации
\[
    \widehat\Delta_i^t
    =
    \frac{\Delta_i^t}{a_i^t},
\]
где \(a_i^t\) описывает эффективную накопленную длину локального
шага. Для обычного SGD с постоянным шагом простейшим ориентиром служит
\(a_i^t=\tau_i\), но при импульсе или меняющихся шагах коэффициент
должен учитывать всю локальную рекурсию. Точная нормализация приводит
к проблеме несогласованности целевой функции и к алгоритму FedNova,
которые рассматриваются в следующих лекциях.

Реальные клиенты различаются также по памяти, доступной мощности,
заряду и качеству соединения. В этой лекции эти факторы служат
мотивацией, но явно не входят в математическую модель. Основные
моделируемые величины --- \(n_i\), \(B_i\), \(E_i\), \(\tau_i\),
\(c_i\), \(\eta_i\) и параметры локального оптимизатора.
